\section{Basic Exercises}

Solve for all least-squares solutions of the equation $A\vx = \vb$
using the normal equations.
\begin{tasks}[style=enumerate](2)
\task
  \begin{displaymath}
    A = \left[\begin{matrix}-3 & 0\\0 & 3\\3 & -6\end{matrix}\right]
    \qquad
    \vb = \left[\begin{matrix}-3\\8\\3\end{matrix}\right]
  \end{displaymath}
\task
  \begin{displaymath}
    A =
    \begin{bmatrix}
      1 & 2 \\
      0 & -1 \\
      1 & -1
    \end{bmatrix}
    \qquad
    \vb = \vThree 4 2 0
  \end{displaymath}
\task
  \begin{displaymath}
    A = \left[\begin{matrix}1 & 1 & -4\\0 & 1 & -3\\-1 & -2 & 7\end{matrix}\right]
    \qquad
    \vb = \left[\begin{matrix}-7\\-10\\2\end{matrix}\right]
  \end{displaymath}
\task
  \begin{displaymath}
    A = \left[\begin{matrix}-1 & 1\\-3 & 1\\-3 & 0\end{matrix}\right]
    \qquad
    \vb = \left[\begin{matrix}1\\11\\15\end{matrix}\right]
  \end{displaymath}
\task
  \begin{displaymath}
    A =
    \begin{bmatrix}
      1 & 0 & 1 \\
      1 & 0 & 1 \\
      1 & 0 & 1 \\
      1 & 1 & 0 \\
      1 & 1 & 0 \\
      1 & 1 & 0 \\
    \end{bmatrix}
    \qquad
    \vb =
    \begin{bmatrix}
      6 \\ 5 \\ 4 \\ 7 \\ 2 \\ 3
    \end{bmatrix}
  \end{displaymath}
\end{tasks}
\begin{tasks}[resume, style=enumerate](1)
\task
  Determine the least squares linear fit to the following data points.
  Sketch the graphs to verify your answer.  You should set up the
  required matrix equations by hand, but you can use a computer to
  determine the coefficients of your model.
  \begin{displaymath}
    \{(-5, 2), (-2, 1), (0, 0), (2, 4), (6, 6)\}
  \end{displaymath}
\task Determine the least squares quadratic fit to the data points
  from the previous problem. Sketch the graphs to verify your answer.
  You should set up the required matrix equations by hand, but you can
  use a computer to determine the coefficients of your model.
\task Determine the design matrix for the following function model
  using the data points from the previous problem.
  \begin{displaymath}
    f(x) = \beta_12^x + \beta_2x\sin(x) + \beta_3x^2 + 7\beta_4
  \end{displaymath}
\task
  Determine the least-squares fit model for the given data.  Round the
  parameters to the nearest hundredth.
  \begin{displaymath}
    \{(-2, 0), (-1, 5), (0, 13), (1, 9), (2, 5), (3, 0)\}
  \end{displaymath}
\task
  Determine the design matrix for the following function model using
  the data points from the previous problem.
  \begin{displaymath}
    f(x) = \beta_1 \cos x + \beta_2 \sin x + \beta_3 x + \beta_4
  \end{displaymath}
\end{tasks}

\section{True/False}

\trueFalseHeader

\begin{tasks}[resume, style=enumerate](1)
\task
  Given any matrix equation $A\vx = \vb$, there always exists a
  least-squares solution.
\task
  Given any matrix equation $A\vx = \vb$, there always exists a unique
  least-squares solution.
\task
  A least-squares solution of $A\vx = \vb$ is a vector $\hat{\vx}$
  that satisfies $A\hat{\vx} = \hat{\vb}$, where $\hat{\vb}$ is the
  orthogonal projection of $\vb$ onto $\col A$.
\task
  Given a model function $f(x) = \beta_0 + \beta_1^2 x$, we may
  minimize the least squares error by using the least squares solution
  to a matrix equation.
\task
  If $X$ denotes the design matrix for a least squares regression
  problem, then $X^T X$ is invertible.
\task
  There are orthogonally diagonalizable matrices that are not symmetric.
\task
  For any square matrix $A$, if $A \vx = \vb$ has a unique least
  squares solution for each choice of $\vb$, then $A$ is invertible.
\task
  Every $n \times n$ symmetric matrix must have $n$ distinct eigenvalues.
\task
  For a symmetric matrix, the dimension of each eigenspace is equal to
  the algebraic multiplicity of the corresponding eigenvalue.
\task
  For any matrix $A \in \R^{m \times n}$ and
  vector $\vb \in \R^m$, if $A\vx = \vb$ has
  a unique least-squares solution, then $\vx \mapsto A\vx$ is onto.
\end{tasks}

\section{More Difficult Problems}

\begin{tasks}[resume, style=enumerate](1)
\task
  Determine a formula for the least-squares solution of $A\vx = \vb$
  given that the columns of $A$ are orthonormal.  Furthermore, explain
  why this solution is unique.
\task
  Let $C$ be an arbitrary $m \times n$ matrix of rank $k$ and let
  $\vb$ be an arbitrary vector in $\R^m$.  Determine the maximum
  size of a linearly independent set of least squares solutions for
  $C\vx = \vb$ where $\vb \not = \vzero$.  Justify your
  answer.  Your solution should be given in terms of $m$, $n$, and
  $k$.
\task
  Consider the following function model:
  \begin{displaymath}
    f(x) = \beta_0 + \beta_1 x^3 + \beta_2 (2^x)
  \end{displaymath}
  where $\beta_0$, $\beta_1$, and $\beta_2$ are parameters.  Determine
  the design matrix which can be used to compute the parameters that
  minimize least squares error for the following data points.
  \begin{displaymath}
    \{(-1, 1), (0, -1), (1, 1), (2, 3), (4, 5)\}
  \end{displaymath}
\task
  \faCalculator \ Consider the following multivariate function model:
  \begin{displaymath}
    f(x,y) = \beta_0 + \beta_1 x + \beta_2 y
  \end{displaymath}
  where $\beta_0$, $\beta_1$, and $\beta_2$ are parameters.  Compute
  the parameters that minimize least squares error for the following
  data points.
  \begin{displaymath}
    \{(-2,0,0), (0,0,3), (0,2,0), (-2, -3, 3), (1,5,-1)\}.
  \end{displaymath}
  Justify your answer.
\end{tasks}

\section{Challenge Problems}

\begin{tasks}[resume, style=enumerate](1)
\task
  Suppose we are given four arbitrary data points:
  \begin{displaymath}
    (x_1, y_1), (x_2, y_2), (x_3, y_3), (x_4, y_4)
  \end{displaymath}

  \begin{enumerate}
  \item
    Construct the design matrix for the given data which can be used
    to find the best-fit curve of the form
    \begin{displaymath}
      f_{\beta_1, \beta_2}(\theta) = \beta_1 \cos\theta + \beta_2
      \sin\theta
    \end{displaymath}
    where $\beta_1$ and $\beta_2$ are parameters.
  \item
    Consider fitting the data with a curve of the form
    \begin{displaymath}
      g_{\alpha_1, \alpha_2}(\theta) = \alpha_1\cos(\theta + \alpha_2)
    \end{displaymath}
    where $\alpha_1$ and $\alpha_2$ are parameters. Note that
    $g_{\alpha_1, \alpha_2}$ is not linear in its parameters.  Given
    $\hat{\alpha_1}$, $\hat{\alpha_2}$, $\hat{\beta_1}$ and
    $\hat{\beta_2}$\textemdash the parameters for the best-fit curves
    \textemdash show that
    \begin{displaymath}
      \sum_{i = 1}^4
      \left(
      \hat{\beta_1}\cos(x_i) + \hat{\beta_2}\sin(x_i) - y_i
      \right)^2
      \leq
      \sum_{i = 1}^4
      \left(
      \hat{\alpha_1} \cos(x_i + \hat{\alpha_2}) - y_i
      \right)^2
    \end{displaymath}
    using the trigonometric identity
    \begin{displaymath}
      \cos(a + b) = \cos(a)\sin(b) - \sin(a)\cos(b)
    \end{displaymath}
    In other words, show that the best-fit curve from part (a) has error
    at least as small as the error of the best-fit curve from part (b).
  \item
    Compute $\hat{\alpha_1}$ and $\hat{\alpha_2}$ from $\hat{\beta_1}$
    and $\hat{\beta_2}$. This implies that, in fact, the errors are
    equal, and that we can fit to the function $g_{\alpha_1,
      \alpha_2}$ using the least-squares method.
  \end{enumerate}
\end{tasks}
